\documentclass[../../main.tex]{subfiles}% 注意这里的文件路径不能用 ./main.tex ,否则用latexmk编译子文件会报错
\graphicspath{{\subfix{./image/}}{\subfix{../../image/}}} % 指定图片引用路径,后续可以直接使用图片文件名
% 如果图片存放在上级目录,则使用 ../../image/ 作为备用路径

% 例如:
% \begin{figure}[H]
% \centering
% \includegraphics[scale=0.3]{图.png}
% \caption{}
% \label{figure:图}
% \end{figure}
% 注意:上述\label{}一定要放在\caption{}之后,否则引用图片序号会只会显示??.

\begin{document}

\section{$\mathscr{S}'$上的Fourier变换}

\begin{definition}
对于$\varphi\in L^1(\mathbb{R}^n)$，$\varphi$的$\mathbf{Fourier}$\textbf{变换}定义为
\begin{align*}
(\mathscr{F}\varphi)(\xi)=\int_{\mathbb{R}^n}\varphi(x)\exp(-2\pi\mathrm{i}x\cdot\xi)\mathrm{d}x,
\end{align*}
其中$x\cdot\xi=x_1\xi_1+x_2\xi_2+\cdots+x_n\xi_n$。

函数$\varphi$的$\mathbf{Fourier}$\textbf{逆变换}定义为
\begin{align*}
(\overline{\mathscr{F}}\varphi)(\xi)=\int_{\mathbb{R}^n}\varphi(x)\exp(2\pi\mathrm{i}x\cdot\xi)\mathrm{d}x.
\end{align*}
\end{definition}

\begin{proposition}
$\mathscr{F}\in\mathscr{L}(\mathscr{S})$.
\end{proposition}
\begin{proof}

注意如下两条性质：
\begin{enumerate}[(1)]
\item $\mathscr{F}(\partial^\alpha\varphi)=(2\pi\mathrm{i}\xi)^\alpha(\mathscr{F}\varphi)(\xi)$,
\item $\mathscr{F}((-2\pi\mathrm{i}x)^\alpha\varphi)(\xi)=\partial^\alpha(\mathscr{F}\varphi)(\xi)$.
\end{enumerate}
于是
\begin{align*}
\|\mathscr{F}\varphi\|_m&=\sup_{\substack{\xi\in\mathbb{R}^n\\|\alpha|\leqslant m}}(1+|\xi|^2)^{\frac{m}{2}}|\partial^\alpha(\mathscr{F}\varphi)(\xi)|\\
&=\sup_{\substack{\xi\in\mathbb{R}^n\\|\alpha|\leqslant m}}\left|\left(\mathscr{F}\left[\left(1-\frac{\Delta}{4\pi^2}\right)^{\frac{m}{2}}(-2\pi\mathrm{i}x)^\alpha\varphi\right]\right)(\xi)\right|\\
&\leqslant\sup_{|\alpha|\leqslant m}\int_{\mathbb{R}^n}\left|\left(1-\frac{\Delta}{4\pi^2}\right)^{\frac{m}{2}}(-2\pi\mathrm{i}x)^\alpha\varphi\right|\mathrm{d}x\\
&\leqslant\sup_{|\alpha|\leqslant m}\sum\limits_{\substack{p\leqslant\alpha\\q\leqslant m}}A_{\alpha,p,q}\int_{\mathbb{R}^n}|x^p\partial^q\varphi|\mathrm{d}x\\
&\leqslant M_m\sup_{\substack{x\in\mathbb{R}^n\\|\alpha|\leqslant m}}(1+|x|^2)^{\frac{m}{2}+n}|\partial^\alpha\varphi(x)|\\
&\leqslant M_m\|\varphi\|_{m+2n}\quad(m=2,4,6,\cdots),
\end{align*}
其中$A_{\alpha,p,q}$及$M_m$均为常数.

\end{proof}

\begin{proposition}
$\overline{\mathscr{F}}=\sigma\mathscr{F}$，从而$\overline{\mathscr{F}}\in\mathscr{L}(\mathscr{S})$.
\end{proposition}

\begin{proposition}
若$\varphi,\psi\in\mathscr{S}$，则
\begin{align}
\langle\mathscr{F}\varphi,\psi\rangle=\langle\varphi,\mathscr{F}\psi\rangle,\label{eq:ftadj1}\\
\langle\overline{\mathscr{F}}\varphi,\psi\rangle=\langle\varphi,\overline{\mathscr{F}}\psi\rangle.\label{eq:ftadj2}
\end{align}
\end{proposition}
\begin{proof}

由Fubini定理，
\begin{align*}
\langle\mathscr{F}\varphi,\psi\rangle&=\int_{\mathbb{R}^n}\left[\int_{\mathbb{R}^n}\varphi(x)\exp(-2\pi\mathrm{i}x\cdot\xi)\mathrm{d}x\right]\psi(\xi)\mathrm{d}\xi\\
&=\int_{\mathbb{R}^n}\varphi(x)\left(\int_{\mathbb{R}^n}\exp(-2\pi\mathrm{i}x\cdot\xi)\psi(\xi)\mathrm{d}\xi\right)\mathrm{d}x\\
&=\langle\varphi,\mathscr{F}\psi\rangle,
\end{align*}
即得\eqref{eq:ftadj1}式。同理可证\eqref{eq:ftadj2}式.

\end{proof}

\begin{definition}
在空间$\mathscr{S}'$上定义$\widetilde{\mathscr{F}}=\mathscr{F}^*$，$\widetilde{\overline{\mathscr{F}}}=(\overline{\mathscr{F}})^*$，称为\textbf{广义}$\mathbf{Fourier}$\textbf{(逆)变换}。在不引起混淆时，可简记$\mathscr{F}\varphi=\widehat{\varphi}$.
\end{definition}

\begin{proposition}
$\mathscr{F}\in\mathscr{L}(\mathscr{S}')$，$\overline{\mathscr{F}}\in\mathscr{L}(\mathscr{S}')$，且限制在$\mathscr{S}$上时，分别与普通Fourier变换、Fourier逆变换一致。
\end{proposition}

\begin{remark}
Fourier变换下微商与乘法、平移与相移的对应关系如下表所示，表中用“$f \circ \text{-----} \cdot g$”表示$g=\mathscr{F}f$或$f=\overline{\mathscr{F}}g$.
\begin{table}[H]
\centering
\caption{}
\begin{tabular}{c|c}
\hline
编号 & 假设 \\
\cline{2-2}
& $f \circ \text{-----} \cdot g$ \\
\hline
(1) & $\widetilde{\partial}^\alpha f \circ \text{-----} \cdot (2\pi\mathrm{i}\xi)^\alpha g$ \\
\hline
(2) & $(-2\pi\mathrm{i}x)^\alpha f \circ \text{-----} \cdot \widetilde{\partial}^\alpha g$ \\
\hline
(3) & $\widetilde{\tau}_a f \circ \text{-----} \cdot \exp(-2\pi\mathrm{i}a\cdot\xi)g$ \\
\hline
(4) & $\exp(2\pi\mathrm{i}a\cdot x)f \circ \text{-----} \cdot \widetilde{\tau}_a g$ \\
\hline
\end{tabular}
\end{table}

\end{remark}

\begin{proposition}
$\mathscr{F}(\exp(-\pi|x|^2))=\exp(-\pi|\xi|^2)$.
\end{proposition}
\begin{proof}

先证$n=1$的情形，令$f(x)=\exp(-\pi x^2)$，则
\begin{align*}
f'(x)+2\pi x f(x)=0,\quad f(0)=1.
\end{align*}
等式两边作Fourier变换，得
\begin{align*}
2\pi\mathrm{i}\xi(\mathscr{F}f)(\xi)+\mathrm{i}(\mathscr{F}f)'(\xi)=0,\quad (\mathscr{F}f)(0)=\int_{-\infty}^{+\infty}\exp(-\pi x^2)\mathrm{d}x=1.
\end{align*}
故$\mathscr{F}f$与$f$满足同一微分方程且初值相同，由常微分方程初值问题解的唯一性，得
\begin{align*}
(\mathscr{F}f)(\xi)=f(\xi)=\exp(-\pi|\xi|^2),\quad \forall\xi\in\mathbb{R}.
\end{align*}
对一般$n\in\mathbb{N}$，利用分离变量即得结论.

\end{proof}

\begin{proposition}
$\mathscr{F}\delta=1$，$\overline{\mathscr{F}}\delta=1$.
\end{proposition}
\begin{proof}

对任意$\varphi\in\mathscr{S}$，
\begin{align*}
\langle\mathscr{F}\delta,\varphi\rangle=\langle\delta,\mathscr{F}\varphi\rangle=(\mathscr{F}\varphi)(0)=\int_{\mathbb{R}^n}\varphi(x)\mathrm{d}x=\langle1,\varphi\rangle.
\end{align*}

\end{proof}

\begin{proposition}
$\mathscr{F}(1)=\delta$，$\overline{\mathscr{F}}(1)=\delta$.
\end{proposition}
\begin{proof}

由$\mathscr{F}\left[\exp\left(-\pi\frac{|x|^2}{m}\right)\right]=m^{\frac{n}{2}}\exp(-m\pi|\xi|^2)$，当$m\to\infty$时，
\begin{align*}
\exp\left(-\pi\frac{|x|^2}{m}\right)\to1\quad(\mathscr{S}'),\quad m^{\frac{n}{2}}\exp(-m\pi|\xi|^2)\to\delta\quad(\mathscr{S}').
\end{align*}
由$\mathscr{F}$的连续性得$\mathscr{F}(1)=\delta$，同理可证另一式.

\end{proof}

\begin{proposition}
\begin{enumerate}[(1)]
\item $\mathscr{F}(p(x))=p\left(\frac{\mathrm{i}}{2\pi}\widetilde{\partial}\right)\delta(\xi)$，其中$p(\cdot)$为多项式；
\item $\mathscr{F}(\widetilde{\partial}^\alpha\delta)=(2\pi\mathrm{i}\xi)^\alpha$.
\end{enumerate}
\end{proposition}
\begin{proof}

(1) 由平移相移对应关系及$\mathscr{F}(1)=\delta$即得；
(2) 由微商乘法对应关系及$\mathscr{F}\delta=1$即得.

\end{proof}

\begin{theorem}
$\overline{\mathscr{F}}=\mathscr{F}^{-1}$，即$\overline{\mathscr{F}}\mathscr{F}=\mathscr{F}\overline{\mathscr{F}}=I$.
\end{theorem}
\begin{proof}

对任意$\varphi\in\mathscr{S}$，$y\in\mathbb{R}^n$，
\begin{align*}
\varphi(y)&=\langle\delta,\tau_{-y}\varphi\rangle=\langle\widetilde{\tau}_y\delta,\varphi\rangle
=\langle\mathscr{F}(\exp(2\pi\mathrm{i}\xi\cdot y)),\varphi\rangle\\
&=\langle\exp(2\pi\mathrm{i}\xi\cdot y),(\mathscr{F}\varphi)(\xi)\rangle
=\int_{\mathbb{R}^n}\exp(2\pi\mathrm{i}\xi\cdot y)(\mathscr{F}\varphi)(\xi)\mathrm{d}\xi=(\overline{\mathscr{F}}\mathscr{F}\varphi)(y),
\end{align*}
故$\varphi=\overline{\mathscr{F}}\mathscr{F}\varphi$，同理$\varphi=\mathscr{F}\overline{\mathscr{F}}\varphi$.

对任意$f\in\mathscr{S}'$，
\begin{align*}
\langle\overline{\mathscr{F}}\mathscr{F}f,\varphi\rangle&=\langle f,\overline{\mathscr{F}}\mathscr{F}\varphi\rangle=\langle f,\varphi\rangle,\quad \forall\varphi\in\mathscr{S},\\
\langle\mathscr{F}\overline{\mathscr{F}}f,\varphi\rangle&=\langle f,\mathscr{F}\overline{\mathscr{F}}\varphi\rangle=\langle f,\varphi\rangle,\quad \forall\varphi\in\mathscr{S},
\end{align*}
因此$\overline{\mathscr{F}}=\mathscr{F}^{-1}$.

\end{proof}

\begin{corollary}
若$f\in L^2$，则$\widetilde{\mathscr{F}}f\in L^2$，且
\begin{align}
\|f\|_{L^2}=\|\widetilde{\mathscr{F}}f\|_{L^2},\label{eq:plancherel}
\end{align}
即$\widetilde{\mathscr{F}}$保持$L^2$范数不变.
\end{corollary}
\begin{remark}
由极化恒等式与上述推论可得：若$f,g\in L^2$，则
\begin{align}
(f,g)_{L^2}=(\widetilde{\mathscr{F}}f,\widetilde{\mathscr{F}}g)_{L^2},\label{eq:innerprodft}
\end{align}
即$\widetilde{\mathscr{F}}$保持$L^2$内积不变.
\end{remark}
\begin{proof}

对任意$\varphi\in\mathscr{S}$，
\begin{align*}
\|\varphi\|_{L^2}^2=\langle\varphi,\overline{\varphi}\rangle=\langle\overline{\mathscr{F}}\mathscr{F}\varphi,\overline{\varphi}\rangle=\langle\mathscr{F}\varphi,\overline{\mathscr{F}}\overline{\varphi}\rangle=\|\mathscr{F}\varphi\|_{L^2}^2.
\end{align*}
因$\mathscr{S}$在$L^2$中稠密，对任意$f\in L^2$，存在$\{\varphi_m\}\subset\mathscr{S}$使得$\|\varphi_m-f\|_{L^2}\to0\ (\lim\limits_{m\to\infty})$.

对任意$p\in\mathbb{N}$，$\|\mathscr{F}\varphi_{m+p}-\mathscr{F}\varphi_m\|_{L^2}\to0\ (\lim\limits_{m\to\infty})$，故$\{\mathscr{F}\varphi_m\}_{m=1}^\infty$是$L^2$中的基本列。结合$\mathscr{F}$在$\mathscr{S}'$上的连续性及$L^2\subseteq\mathscr{S}'$，得
\begin{align*}
\widetilde{\mathscr{F}}f=\lim\limits_{m\to\infty}\mathscr{F}\varphi_m\quad(L^2),
\end{align*}
因此$\widetilde{\mathscr{F}}f\in L^2$，且
\begin{align*}
\|\widetilde{\mathscr{F}}f\|_{L^2}=\lim\limits_{m\to\infty}\|\mathscr{F}\varphi_m\|_{L^2}=\lim\limits_{m\to\infty}\|\varphi_m\|_{L^2}=\|f\|_{L^2}.
\end{align*}

\end{proof}



\begin{example}
设$H^m(\mathbb{R}^n) = \{u \in \mathscr{S}' \mid \widetilde{\partial}^\alpha u \in L^2(\mathbb{R}^n)\ (|\alpha| \leqslant m)\}$，其中范数定义为
\begin{align*}
\|u\|_m = \left(\sum\limits_{|\alpha|\leqslant m}\|\widetilde{\partial}^\alpha u\|_{L^2}^2\right)^{\frac12}.
\end{align*}
又对$\forall u \in H^m(\mathbb{R}^n)$，定义
\begin{align*}
\|u\|_m' = \left(\int_{\mathbb{R}^n}(1+|\xi|^2)^m|(\mathscr{F}u)(\xi)|^2\mathrm{d}\xi\right)^{\frac12},
\end{align*}
求证:
\begin{enumerate}[(1)]
\item $\|u\|_m' < \infty$;
\item $\|\cdot\|_m'$是$H^m(\mathbb{R}^n)$的等价范数;
\item $H^m(\mathbb{R}^n)$是完备的.
\end{enumerate}
\end{example}

\begin{example}
对任意的非负实数$s$，设
\begin{align*}
H^s(\mathbb{R}^n) \triangleq \{u \in L^2(\mathbb{R}^n) \mid (1+|\xi|^2)^{s/2}\widehat{u}(\xi) \in L^2(\mathbb{R}^n)\},
\end{align*}
其中范数定义为
\begin{align*}
\|u\|_s = \left\|(1+|\xi|^2)^{s/2}\widehat{u}(\xi)\right\|_{L^2}.
\end{align*}
求证:
\begin{enumerate}[(1)]
\item 当$s = m \in \mathbb{N}$时，这种定义与原来$H^m(\mathbb{R}^n)$的定义等价;
\item $H^s(\mathbb{R}^n)$中可引进内积$(\cdot,\cdot)$，使得$\|u\|_s = (u,u)^{\frac12}$;
\item 设$u \in H^s(\mathbb{R}^n)'$，求证: 存在$\widetilde{u} \in L^1_{\mathrm{loc}}(\mathbb{R}^n)$，使得
\begin{align*}
\widetilde{u}(\xi)(1+|\xi|^2)^{-s/2} \in L^2(\mathbb{R}^n),
\end{align*}
并且
\begin{align*}
\langle u,\mathscr{F}\varphi\rangle = \int_{\mathbb{R}^n}\varphi(\xi)\cdot\widetilde{u}(\xi)\mathrm{d}\xi \quad (\forall \varphi \in \mathscr{S}).
\end{align*}
\end{enumerate}
\end{example}

\begin{example}
设$f(x) \in L^1(\mathbb{R}^n)$，求证:
\begin{align*}
(\mathscr{F}f)(\xi) = \int_{\mathbb{R}^n}f(x)e^{-2\pi i x\cdot\xi}\mathrm{d}x,
\end{align*}
即$f(x)$按$\mathscr{S}'$的Fourier变换与普通的Fourier变换一致.
\end{example}

\begin{example}
求证: 方程$\Delta f = f$在$\mathscr{S}'(\mathbb{R}^n)$中无非零解.
\end{example}












\end{document}